% Generated draft: requires mathematical and source review.
\documentclass[11pt]{article}
\usepackage[margin=25mm]{geometry}
\usepackage{fontspec}
\setmainfont{Latin Modern Roman}
\usepackage{amsmath,amssymb,mathtools,unicode-math}
\setmathfont{Latin Modern Math}
\usepackage{enumitem,needspace,xurl,hyperref}
\setlength{\parindent}{0pt}
\setlength{\parskip}{5pt}
\setlength{\emergencystretch}{3em}
\newcommand{\N}{\mathbb{N}}
\newcommand{\Z}{\mathbb{Z}}
\newcommand{\Q}{\mathbb{Q}}
\newcommand{\R}{\mathbb{R}}
\newcommand{\C}{\mathbb{C}}
\newcommand{\F}{\mathbb{F}}
\newcommand{\A}{\mathbb{A}}
\newcommand{\PP}{\mathbb{P}}
\newcommand{\E}{\mathbb{E}}
\newcommand{\eps}{\varepsilon}
\newcommand{\dd}{\,\mathrm{d}}
\begin{document}
\section*{UnsolvedMath ID: 30006399}
\textbf{Source problem code:} catalogue entry
\subsection{Definitions and mathematical statement}
The recoverable displayed expression is
\[
f_n(z)=\sum_{j=0}^{n}\xi_j z^j.
\]
For purposes of interpreting this expression, take \(n\) to be a nonnegative integer and \(z\) a complex variable. The symbols \(\xi_j\) denote random coefficients, but the supplied statement does not specify their values, marginal distributions, dependence, or probability space. These omissions prevent recovery of a particular random-polynomial ensemble.

For any nonzero realization with complex coefficients, its roots can be represented by the multiset
\[
\mathcal Z_n=\{z\in\C:f_n(z)=0\},
\]
where roots are counted with algebraic multiplicity. Its cardinality equals the degree of that realization, which equals \(n\) only if \(\xi_n\ne0\). The supplied statement does not exclude the zero polynomial or specify how such a realization should be treated.

The source question is: \emph{What can be said about the distribution of the roots of \(f_n\)?} It does not specify whether this means the joint law of \(\mathcal Z_n\), an expected or empirical root measure, an asymptotic distribution as \(n\to\infty\), or another root statistic. Although the title and background identify Littlewood polynomials, they do not recover the missing coefficient assumptions or a precise target assertion.

\textbf{Unresolved source qualification.} The record marks the statement as \texttt{unrecoverable}. Its supplied verification says that the root-distribution question is a lead-in to theorems proved in the reported work, rather than a separately posed open problem, and that no clean statement is promoted. Consequently, this transcription does not assert a recovered conjecture or a specific unsolved mathematical claim.
\subsection{Short English statement}
The record asks broadly about the roots of a random polynomial, but omits the coefficient law and the intended meaning of root distribution. Its supplied verification identifies the question as an introductory lead-in to reported theorems. A precise standalone open problem cannot be recovered from this evidence.
\subsection{Sources}
\begin{itemize}
\item Supplied local UnsolvedMath registry record, permanent ID \(30006399\), exact source problem code \texttt{catalogue entry}; collection: Oberwolfach Reports Open Problems. The record lists proposed year \(2025\) and supplies no proposer attribution.
\item Supplied source citation: \emph{Combinatorics, Probability and Computing (2026)}, Oberwolfach Reports, DOI \texttt{10.4171/owr/2025/42}, pp.\ 2243--2322. \url{https://doi.org/10.4171/owr/2025/42}. This is supplied report provenance; the complete report text was not supplied or independently consulted here.
\item Supplied literature source: \emph{Oberwolfach Report 42 2025}. \url{https://ems.press/content/serial-article-files/52246}.
\item Supplied literature source: \emph{Limit Law for Root Separation in Random Polynomials}. \url{https://arxiv.org/abs/2505.02723}.
\item Literature assessment checked at \texttt{2026-08-22}: supplied dated curation, not independently verified. It labels the broad question open and reports that bulk root location, angular equidistribution, and recently a minimum-spacing law are known for the Littlewood ensemble, while a broad request for the full root distribution remains open rather than constituting one settled theorem. The background explicitly calls for independent verification of time-sensitive status and cited claims.
\item Supplied statement-verification qualification: the official report uses this question as a lead-in to proved theorems, not as a separately posed open problem. The record retains the original wording for provenance, marks its statement \texttt{unrecoverable}, and nevertheless carries an \texttt{open} status. Its title is \emph{Resolved Root-Distribution Lead-In for Littlewood Polynomials}. These supplied labels do not independently establish mathematical status.
\end{itemize}
\end{document}
